\paragraph{H-A: Net section x load-path straightness controls strength} \textbf{Observation}: At matched porosity 0.20 the strength varies 6x (slit\_90 4.0, ellipse\_0 8.6, random Voronoi 8.2-9.7, fine meshes 11.7-13.3, coarse meshes 16.4-17.0, ellipse\_90 19.0, slit\_0 25.4 N/m). Dividing by the minimum load-bearing solid fraction across the loading direction (minSF) collapses the range to a 'net-section strength' of 12-32 N/m that is highest for straight, uniform-width ligaments parallel to the load (slit\_0 31.8, square mesh 31, p32 22) and lowest for necked/curved paths (round-pore meshes 17-19, random Voronoi 12). \textbf{Proposed mechanism}: sigma\_max = minSF x sigma\_0 / K, where sigma\_0 ~ 39 N/m is the intrinsic strength and K is a geometric knockdown set by the curvature/necking of the load-bearing ligaments (K ~ 1.2 straight strips, ~2 necked round-pore ligaments, ~3 random networks); hierarchy and porosity act only through minSF and K. \textbf{Competing explanation}: H-B: strength is controlled by the fraction of under-coordinated edge atoms (atomic-scale edge weakness), not by continuum stress concentration; at fixed minSF, strength decreases with edge-atom fraction and is insensitive to ligament shape. \textbf{Predicted if mechanism holds}: A coarse square-pore mesh and a coarse round-pore mesh with the same minSF (pore width = square side) have square $>$= round (straight strips beat necked paths) despite more edge atoms in the square; strength of any new design is predicted within ~20\% by minSF x 31/K. \textbf{Predicted if competing explanation holds}: Round $>$= square at equal minSF (fewer edge atoms), and designs with equal edge-atom fraction but different ligament straightness (slit\_0 vs jittered mesh, both fUC ~0.15) would have similar strength -- they differ 2x, so H-B is already disfavoured. \textbf{Discriminating experiments}: S4 pair 1: square vs circle coarse mesh at equal minSF and equal period; S4 pair 2: same at equal porosity; S5: vein-width sweep of nested meshes (rule of mixtures of straight veins and necked fine ligaments). \textbf{Outcome}: PARTLY REFUTED (Stage 4 pairs 1-2): at equal net section, circular coarse pores are as strong as square ones (16.1 vs 17.0 N/m; predicted 11.3), and at equal porosity 14.1 (predicted 9.8): the pore shape does not matter at the coarse scale, so the 'necking knockdown' of H-A is wrong. The net-section factor itself is confirmed (slit\_0 at porosity 0.1 and 0.2: 25.6 vs 25.4 N/m; vein-width series). REVISED (H-A'): sigma = minSF x sigma\_net, where sigma\_net depends on the absolute ligament width (edge/discreteness effect, ~17 N/m for 10-13 A ligaments, 22-31 N/m for $>$= 18 A) and on load-path alignment (random networks 12-15 N/m), not on pore shape. Edge-atom fraction (H-B) explains the size dependence but not the orientation effects.

\paragraph{H-C: Fracture mode is set by the equivalence of parallel load paths} \textbf{Observation}: Perfectly periodic straight-ligament structures fail in one avalanche (square mesh: 6 harmless corner events then all rows snap at 17.0 N/m; slit\_0: stepwise with 3 events; clustered vacancies), whereas structures whose load-bearing paths are non-equivalent (fine round-pore meshes 15-22 events, jittered meshes 17-22, nested meshes 15-16, distributed vacancies 9) fail progressively with post-peak load retention. \textbf{Proposed mechanism}: When all load-bearing ligaments are crystallographically equivalent they reach the instability at the same strain and snap together (no load transfer possible); necking, disorder or hierarchy make ligaments non-equivalent so they fail sequentially and the surviving ones carry the transferred load. \textbf{Competing explanation}: H-C': the mode is set by the ligament shape alone (straight = abrupt, necked = progressive) irrespective of equivalence. \textbf{Predicted if mechanism holds}: Introducing +-15\% dispersion of the ligament widths into the square coarse mesh (same porosity) converts the single avalanche into a multi-step failure at nearly the same peak stress; a perfectly uniform fine mesh remains progressive only because of necking. \textbf{Predicted if competing explanation holds}: The dispersed square mesh still fails in one avalanche. \textbf{Discriminating experiments}: S4 pair 3: square mesh with size dispersion 0.15 vs ordered; S4 pair 4: p32 round mesh with and without dispersion. \textbf{Outcome}: REFUTED (Stage 4 pairs 3-4): +-15\% ligament dispersion lowered the strength of both coarse meshes (17.0 -$>$ 15.8; 16.4 -$>$ 13.8) and did not make failure progressive (square: still a single avalanche; p32: shorter, more localised failure, L 0.12 -$>$ 0.46). Load-path non-equivalence is not sufficient; a fibre-bundle-type criterion applies: sequential failure requires enough parallel paths that the load transferred by one failing path stays below the strength margin of the others. Fine meshes (8-10 rows) and jittered meshes satisfy it, three wide strips do not.

\paragraph{H-D: Disorder trades strength for progressivity with an optimum for energy absorption} \textbf{Observation}: Pore-position jitter 0/1/2/3 A gives 11.7/12.7/11.8/11.1 N/m but failure strains 0.13/0.21/0.24/0.22 and 6/17/22/21 damage steps; Voronoi regularity 0/0.6/1 gives 9.7/14.4/13.3 N/m and W 0.74/1.56/1.20 J/m2 (n25); random n50 8.2 N/m. \textbf{Proposed mechanism}: Disorder acts through two opposing channels: it creates weak links (thin or misoriented ligaments) that lower the peak (extreme-value control), and it breaks the equivalence of load paths so that failure becomes sequential (longer tail, more work). The net effect on work to failure is positive at moderate disorder and negative when weak links dominate. \textbf{Competing explanation}: H-D': disorder always lowers both strength and work (pure weak-link picture); or disorder always raises work (pure crack-path-disruption picture). \textbf{Predicted if mechanism holds}: W(disorder) is non-monotonic with a maximum at intermediate disorder for both jittered meshes and Voronoi networks; seed-to-seed scatter of strength grows with disorder. \textbf{Predicted if competing explanation holds}: W monotonically decreasing (weak link) or increasing (disruption). \textbf{Discriminating experiments}: S5: Voronoi regularity sweep 0/0.3/0.6/0.8/1.0; S6: 3 seeds each for jitter 2 A, size dispersion 0.15, Voronoi reg 0.0 and 0.6; lattice-registry replicates of the ordered mesh. \textbf{Outcome}: PARTLY SUPPORTED, seed statistics decisive (Stages 5-6): jitter trades a mild strength loss (12.7 -$>$ 11.1 N/m) for longer tails, but seed replicates show large scatter of the post-peak tail (jitter 2 A: failure strain 0.24 vs 0.15 between seeds). In Voronoi networks the apparent optimum at regularity 0.6 (14.4 N/m) was a single-seed effect (seed 2: 9.2 N/m); regularity 0.3/0.8 give 10.5/6.6 N/m and random seeds 8.7-9.7 N/m: strength of random networks is controlled by the weakest ligament (extreme-value statistics), not by regularity, and work does not show a robust optimum. Lattice-registry replicates of the ordered mesh (11.7/12.0/13.1/13.1/13.3 N/m) show that the atomic placement of pore edges alone moves strength by ~12\% and can switch the mode (abrupt vs progressive), setting the resolution of all fine-mesh comparisons.

\paragraph{H-E: Flaw tolerance beyond ~20 A: crack-tip lattice trapping caps the size effect} \textbf{Observation}: Cracks of 10/20/30 A give 24.1/18.5/18.7 N/m (Griffith scaling would give 24.1/17.0/13.9); the 30 A crack starts breaking bonds at 7.9\% strain and grows stably to 9.9\% before running; a 20 A hole (19.1) equals a 20 A crack. \textbf{Proposed mechanism}: In this discrete lattice the crack advances by discrete bond-breaking events with an energy barrier (lattice trapping); once the flaw exceeds a few nanometres the peak stress is set by the trapping strength of the tip bond configuration rather than by the flaw length. \textbf{Competing explanation}: H-E': the plateau is a periodic-image artefact (30 A crack in a 100 A cell interacts with its images) and strength keeps decreasing with crack length in larger cells. \textbf{Predicted if mechanism holds}: A 40 A crack in a 150 A cell and a 30 A crack in a 150 A cell both give 17-19 N/m; a 30 A hole gives ~18 N/m. \textbf{Predicted if competing explanation holds}: Strength continues to fall (40 A: ~13 N/m in the large cell). \textbf{Discriminating experiments}: S4: precrack 30/40 A in 150 A cells; hole 30 A. \textbf{Outcome}: SUPPORTED (Stage 4): 30 A and 40 A cracks in 150 A cells give 19.5 and 17.9 N/m (predicted 17-19; Griffith scaling 13.9 and 12.0), a 30 A hole 18.6 N/m (20 A hole 19.1). Beyond ~20 A the strength of a flawed sheet is set by the trapping strength of the tip bond configuration (stable bond-by-bond growth precedes the instability), not by the flaw length; the plateau is not a periodic-image artefact.

\paragraph{H-F: Nested hierarchy does not add strength; the finest load-bearing level and the load-parallel veins control it} \textbf{Observation}: Two-level meshes (12.7-14.3) lie between the fine (11.7-13.3) and coarse (16.4-17.0) single-scale limits at matched porosity; strength rises with vein width (W4/8/12: 13.2/12.8/13.7) and falls with domain size (D24/40/60: 14.3/12.8/12.7); the three-level mesh (13.0) equals the fine mesh. Fracture panels show the fine level failing first at pore edges, load-parallel veins carrying the residual load, and cracks channelling along columns next to load-perpendicular veins. \textbf{Proposed mechanism}: The nested structure is a parallel arrangement of straight veins (strong, straight strips) and a necked fine mesh (weak); the fine level reaches its instability first, the veins then carry a fraction W/D of the net section; perpendicular veins act as compliant crack channels rather than crack arrestors. \textbf{Competing explanation}: H-F': veins arrest cracks (compartmentalisation) so that nested meshes gain failure strain and work relative to single-scale meshes of equal strength. \textbf{Predicted if mechanism holds}: Veins along the load only (vein\_dirs='x') give higher strength than veins in both directions at equal porosity, and veins perpendicular only ('y') give no more than the fine mesh; the work to failure of nested meshes does not exceed that of the fine mesh with the same porosity. \textbf{Predicted if competing explanation holds}: xy veins beat x-only veins (compartments) and nested meshes have larger work than single-scale meshes. \textbf{Discriminating experiments}: S5: vein direction x / y / xy at matched porosity; vein width sweep to W = 20 A; hierarchy x disorder and hierarchy x anisotropy interactions. \textbf{Outcome}: SUPPORTED and sharpened (Stage 4-5): veins only along the load give 16.1 N/m and the largest work of the nested family (1.6 J/m2), veins in both directions 12.8, veins only perpendicular 12.0 = fine mesh; the vein-width series 4/8/12/16/20 A (13.2/12.8/13.7/14.8/15.6 N/m) approaches the coarse-only limit (17.0) monotonically and the mode turns abrupt at 20 A. Hierarchy adds strength only through load-parallel coarse ligaments; perpendicular veins are dead weight and act as crack channels rather than crack arrestors (H-F' rejected). Stage 5 additions: scale ratio at fixed vein geometry does not matter (period-8 fine level 12.7 vs period-12 12.8 N/m); D60/W12 13.0; jittering the fine level costs 2.5 N/m (12.8 -$>$ 10.3); elongating the fine pores along the load inside the nested mesh gives 17.5 N/m with W = 2.0 J/m2 and progressive failure -- the only interaction that beats the coarse-only limit (17.0, abrupt), because it combines straight fine ligaments (high net-section strength) with many parallel paths (fibre-bundle progressivity).